% Options for packages loaded elsewhere
% Options for packages loaded elsewhere
\PassOptionsToPackage{unicode}{hyperref}
\PassOptionsToPackage{hyphens}{url}
%
\documentclass[
  14pt,
  ignorenonframetext,
  aspectratio=169,
]{beamer}
\newif\ifbibliography
\usepackage{pgfpages}
\setbeamertemplate{caption}[numbered]
\setbeamertemplate{caption label separator}{: }
\setbeamercolor{caption name}{fg=normal text.fg}
\beamertemplatenavigationsymbolsempty
% remove section numbering
\setbeamertemplate{part page}{
  \centering
  \begin{beamercolorbox}[sep=16pt,center]{part title}
    \usebeamerfont{part title}\insertpart\par
  \end{beamercolorbox}
}
\setbeamertemplate{section page}{
  \centering
  \begin{beamercolorbox}[sep=12pt,center]{section title}
    \usebeamerfont{section title}\insertsection\par
  \end{beamercolorbox}
}
\setbeamertemplate{subsection page}{
  \centering
  \begin{beamercolorbox}[sep=8pt,center]{subsection title}
    \usebeamerfont{subsection title}\insertsubsection\par
  \end{beamercolorbox}
}
% Prevent slide breaks in the middle of a paragraph
\widowpenalties 1 10000
\raggedbottom
\usepackage{iftex}
\ifPDFTeX
  \usepackage[T1]{fontenc}
  \usepackage[utf8]{inputenc}
  \usepackage{textcomp} % provide euro and other symbols
\else % if luatex or xetex
  \usepackage{unicode-math} % this also loads fontspec
  \defaultfontfeatures{Scale=MatchLowercase}
  \defaultfontfeatures[\rmfamily]{Ligatures=TeX,Scale=1}
\fi
\usepackage{lmodern}

\usetheme[]{monash}
\ifPDFTeX\else
  % xetex/luatex font selection
\fi
% Use upquote if available, for straight quotes in verbatim environments
\IfFileExists{upquote.sty}{\usepackage{upquote}}{}
\IfFileExists{microtype.sty}{% use microtype if available
  \usepackage[]{microtype}
  \UseMicrotypeSet[protrusion]{basicmath} % disable protrusion for tt fonts
}{}
\makeatletter
\@ifundefined{KOMAClassName}{% if non-KOMA class
  \IfFileExists{parskip.sty}{%
    \usepackage{parskip}
  }{% else
    \setlength{\parindent}{0pt}
    \setlength{\parskip}{6pt plus 2pt minus 1pt}}
}{% if KOMA class
  \KOMAoptions{parskip=half}}
\makeatother

\usepackage{color}
\usepackage{fancyvrb}
\newcommand{\VerbBar}{|}
\newcommand{\VERB}{\Verb[commandchars=\\\{\}]}
\DefineVerbatimEnvironment{Highlighting}{Verbatim}{commandchars=\\\{\}}
% Add ',fontsize=\small' for more characters per line
\usepackage{framed}
\definecolor{shadecolor}{RGB}{248,248,248}
\newenvironment{Shaded}{\begin{snugshade}}{\end{snugshade}}
\newcommand{\AlertTok}[1]{\textcolor[rgb]{0.94,0.16,0.16}{#1}}
\newcommand{\AnnotationTok}[1]{\textcolor[rgb]{0.56,0.35,0.01}{\textbf{\textit{#1}}}}
\newcommand{\AttributeTok}[1]{\textcolor[rgb]{0.13,0.29,0.53}{#1}}
\newcommand{\BaseNTok}[1]{\textcolor[rgb]{0.00,0.00,0.81}{#1}}
\newcommand{\BuiltInTok}[1]{#1}
\newcommand{\CharTok}[1]{\textcolor[rgb]{0.31,0.60,0.02}{#1}}
\newcommand{\CommentTok}[1]{\textcolor[rgb]{0.56,0.35,0.01}{\textit{#1}}}
\newcommand{\CommentVarTok}[1]{\textcolor[rgb]{0.56,0.35,0.01}{\textbf{\textit{#1}}}}
\newcommand{\ConstantTok}[1]{\textcolor[rgb]{0.56,0.35,0.01}{#1}}
\newcommand{\ControlFlowTok}[1]{\textcolor[rgb]{0.13,0.29,0.53}{\textbf{#1}}}
\newcommand{\DataTypeTok}[1]{\textcolor[rgb]{0.13,0.29,0.53}{#1}}
\newcommand{\DecValTok}[1]{\textcolor[rgb]{0.00,0.00,0.81}{#1}}
\newcommand{\DocumentationTok}[1]{\textcolor[rgb]{0.56,0.35,0.01}{\textbf{\textit{#1}}}}
\newcommand{\ErrorTok}[1]{\textcolor[rgb]{0.64,0.00,0.00}{\textbf{#1}}}
\newcommand{\ExtensionTok}[1]{#1}
\newcommand{\FloatTok}[1]{\textcolor[rgb]{0.00,0.00,0.81}{#1}}
\newcommand{\FunctionTok}[1]{\textcolor[rgb]{0.13,0.29,0.53}{\textbf{#1}}}
\newcommand{\ImportTok}[1]{#1}
\newcommand{\InformationTok}[1]{\textcolor[rgb]{0.56,0.35,0.01}{\textbf{\textit{#1}}}}
\newcommand{\KeywordTok}[1]{\textcolor[rgb]{0.13,0.29,0.53}{\textbf{#1}}}
\newcommand{\NormalTok}[1]{#1}
\newcommand{\OperatorTok}[1]{\textcolor[rgb]{0.81,0.36,0.00}{\textbf{#1}}}
\newcommand{\OtherTok}[1]{\textcolor[rgb]{0.56,0.35,0.01}{#1}}
\newcommand{\PreprocessorTok}[1]{\textcolor[rgb]{0.56,0.35,0.01}{\textit{#1}}}
\newcommand{\RegionMarkerTok}[1]{#1}
\newcommand{\SpecialCharTok}[1]{\textcolor[rgb]{0.81,0.36,0.00}{\textbf{#1}}}
\newcommand{\SpecialStringTok}[1]{\textcolor[rgb]{0.31,0.60,0.02}{#1}}
\newcommand{\StringTok}[1]{\textcolor[rgb]{0.31,0.60,0.02}{#1}}
\newcommand{\VariableTok}[1]{\textcolor[rgb]{0.00,0.00,0.00}{#1}}
\newcommand{\VerbatimStringTok}[1]{\textcolor[rgb]{0.31,0.60,0.02}{#1}}
\newcommand{\WarningTok}[1]{\textcolor[rgb]{0.56,0.35,0.01}{\textbf{\textit{#1}}}}

\usepackage{longtable,booktabs,array}
\newcounter{none} % for unnumbered tables
\usepackage{calc} % for calculating minipage widths
\usepackage{caption}
% Make caption package work with longtable
\makeatletter
\def\fnum@table{\tablename~\thetable}
\makeatother
\usepackage{graphicx}
\makeatletter
\newsavebox\pandoc@box
\newcommand*\pandocbounded[1]{% scales image to fit in text height/width
  \sbox\pandoc@box{#1}%
  \Gscale@div\@tempa{\textheight}{\dimexpr\ht\pandoc@box+\dp\pandoc@box\relax}%
  \Gscale@div\@tempb{\linewidth}{\wd\pandoc@box}%
  \ifdim\@tempb\p@<\@tempa\p@\let\@tempa\@tempb\fi% select the smaller of both
  \ifdim\@tempa\p@<\p@\scalebox{\@tempa}{\usebox\pandoc@box}%
  \else\usebox{\pandoc@box}%
  \fi%
}
% Set default figure placement to htbp
\def\fps@figure{htbp}
\makeatother





\setlength{\emergencystretch}{3em} % prevent overfull lines

\providecommand{\tightlist}{%
  \setlength{\itemsep}{0pt}\setlength{\parskip}{0pt}}



 


% Colors
\definecolor{shadecolor}{RGB}{225,225,225}
\setbeamercolor{description item}{fg=Orange}
\definecolor{MonashBlue}{RGB}{66,109,152}
\definecolor{burntorange}{rgb}{0.8, 0.33, 0.0}

% Packages
\usepackage{amsmath, bm, amssymb, amsthm, mathrsfs,pifont}
\usepackage{url}
\usepackage{multirow, booktabs, float, textcmds, siunitx}
\usepackage{bm,booktabs,animate,ragged2e,multicol,microtype,hyperref}
\usepackage{fontawesome5}

% Figures
\graphicspath{{figs/}}

% Fonts
\fontsize{13}{15}\sf
\setbeamerfont{title}{series=\bfseries,parent=structure,size=\fontsize{24}{24}}
\ifcsname Shaded\endcsname
  \definecolor{shadecolor}{RGB}{225,225,225}
  \renewenvironment{Shaded}{\vspace*{0.15cm}\color{black}\fontsize{10}{10}\sf\begin{snugshade}\color{black}}{\end{snugshade}}
\fi

% gt dependencies
\usepackage{longtable,caption,setspace}
\captionsetup{font={small,stretch=0.80}}

% My definitions

\def\E{\text{E}}
\def\V{\text{Var}}
\def\up#1{\raisebox{-0.3cm}{#1}}
\def\pred#1#2#3{\hat{#1}_{#2|#3}}
\def\damped{$_\text{d}$}
\def\h+{h_{m}^{+}}
\def\str#1{\rlap{#1}\textcolor{red}{\rule{1cm}{0.1cm}}}

\def\st#1{\rlap{#1}\textcolor{red}{\rule{1cm}{0.1cm}}}
\def\bY{\bm{y}}
\def\by{\bm{y}}
\def\bS{\bm{S}}
\def\bI{\text{\rm\textbf{I}}}
\def\bbeta{\bm{\beta}}
\def\bSigma{\bm{\Sigma}}
\def\bW{\bm{\Sigma}}
\def\Var{\text{Var}}
\def\var{\text{Var}}
\def\bOmega{\bm{\Omega}}
\def\bLambda{\bm{\Lambda}}
\let\mc\multicolumn
\def\hl{\color[RGB]{230, 172, 0}}

\def\forecast{\begin{alertblock}{}\fontsize{10}{11}\sf
 A forecast is an estimate of the probability distribution of a variable to be observed in the future.
\end{alertblock}}

\def\simfutures{\begin{textblock}{2.9}(12.8,7.7)
\begin{block}{}\fontsize{10}{11}\sf
Simulated futures from an ETS model
\end{block}\end{textblock}}

\setbeamertemplate{title page}
{\placefig{0}{-0.5}{width=1.01\paperwidth,height=1.45\paperheight}{africa1.jpg}
\begin{textblock}{8.5}(.5,0.2)
  \raggedright\fontsize{20}{20}\selectfont\bfseries\sffamily\textcolor{white}{\inserttitle}\\[0.2cm]
  \fontsize{12}{12}\selectfont\bfseries\sffamily\textcolor{white}{\insertsubtitle}
  \end{textblock}
\placefig{.5}{6.2}{width=2.4cm}{tsibble.png}
\placefig{2.9}{6.2}{width=2.4cm}{feasts.png}
\placefig{5.3}{6.2}{width=2.4cm}{fable.png}
\begin{textblock}{7.5}(10.6,-.1)
  {\fontsize{9}{9}\sf\color{white}https://workshop.f4sg.org/africast/}
\end{textblock}
}

% Tikz plots
\usepackage{tikz}
\usetikzlibrary{trees,shapes,arrows,matrix,shadows,positioning,calc}
\tikzstyle{decision} = [diamond, draw, fill=blue!20,
    text width=4.5em, text badly centered, node distance=4cm, inner sep=0pt]
\tikzstyle{block} = [rectangle, draw, fill=blue!20,
    text width=5cm, text centered, rounded corners, minimum height=4em]
\tikzstyle{line} = [draw, thick, -latex']
\tikzstyle{cloud} = [draw, ellipse,fill=red!20, node distance=3cm,
    minimum height=2em, text centered]
\tikzstyle{connector} = [->,thick]

\tikzset{
  basic/.style  = {draw, text width=2cm, font=\sffamily, rectangle},
  root/.style   = {basic, text width=3cm, rounded corners=2pt, thin, align=center, fill=red!30},
  level 2a/.style = {basic, rounded corners=2pt, thin,align=center, fill=blue!50, text width=7em},
  level 2b/.style = {basic, rounded corners=2pt, thin,align=center, fill=green!50, text width=7em},
  level 3a/.style = {basic, rounded corners=2pt, thin, align=center, fill=blue!30, text width=4em},
  level 3b/.style = {basic, rounded corners=2pt, thin, align=center, fill=green!30, text width=4em},
  level 4a/.style = {basic, rounded corners=2pt, thin, align=left, fill=blue!10, text width=3.5em},
  level 4b/.style = {basic, rounded corners=2pt, thin, align=left, fill=green!10, text width=3.5em}
}

% % BIBLIOGRAPHIES
% \usepackage[style=authoryear,bibencoding=utf8,minnames=1,maxnames=4, maxbibnames=99,natbib=true,dashed=false,terseinits=true,giveninits=true,uniquename=false,uniquelist=false,labeldate=true,doi=false, isbn=false, natbib=true,backend=biber]{biblatex}

% \DeclareFieldFormat{url}{\url{#1}}
% \DeclareFieldFormat[article]{pages}{#1}
% \DeclareFieldFormat[inproceedings]{pages}{\lowercase{pp.}#1}
% \DeclareFieldFormat[incollection]{pages}{\lowercase{pp.}#1}
% \DeclareFieldFormat[article]{volume}{\mkbibbold{#1}}
% \DeclareFieldFormat[article]{number}{\mkbibparens{#1}}
% \DeclareFieldFormat[article]{title}{\MakeCapital{#1}}
% \DeclareFieldFormat[article]{url}{}
% \DeclareFieldFormat[Techreport]{Url}{}
% \DeclareFieldFormat[book]{url}{}
% \DeclareFieldFormat[inbook]{url}{}
% \DeclareFieldFormat[incollection]{url}{}
% \DeclareFieldFormat[inproceedings]{url}{}
% \DeclareFieldFormat[inproceedings]{title}{#1}
% \DeclareFieldFormat{shorthandwidth}{#1}
% %\DeclareFieldFormat{extrayear}{}
% % No dot before number of articles
% \usepackage{xpatch}
% \xpatchbibmacro{volume+number+eid}{\setunit*{\adddot}}{}{}{}
% % Remove In: for an article.
% \renewbibmacro{in:}{%
%   \ifentrytype{article}{}{%
%   \printtext{\bibstring{in}\intitlepunct}}}

% \AtEveryBibitem{\clearfield{month}}
% \AtEveryCitekey{\clearfield{month}}
% \AtBeginBibliography{\fontsize{11}{11}\sf}
\setbeamertemplate{frametitle continuation}{}
\makeatletter
\@ifpackageloaded{caption}{}{\usepackage{caption}}
\AtBeginDocument{%
\ifdefined\contentsname
  \renewcommand*\contentsname{Table of contents}
\else
  \newcommand\contentsname{Table of contents}
\fi
\ifdefined\listfigurename
  \renewcommand*\listfigurename{List of Figures}
\else
  \newcommand\listfigurename{List of Figures}
\fi
\ifdefined\listtablename
  \renewcommand*\listtablename{List of Tables}
\else
  \newcommand\listtablename{List of Tables}
\fi
\ifdefined\figurename
  \renewcommand*\figurename{Figure}
\else
  \newcommand\figurename{Figure}
\fi
\ifdefined\tablename
  \renewcommand*\tablename{Table}
\else
  \newcommand\tablename{Table}
\fi
}
\@ifpackageloaded{float}{}{\usepackage{float}}
\floatstyle{ruled}
\@ifundefined{c@chapter}{\newfloat{codelisting}{h}{lop}}{\newfloat{codelisting}{h}{lop}[chapter]}
\floatname{codelisting}{Listing}
\newcommand*\listoflistings{\listof{codelisting}{List of Listings}}
\makeatother
\makeatletter
\makeatother
\makeatletter
\@ifpackageloaded{caption}{}{\usepackage{caption}}
\@ifpackageloaded{subcaption}{}{\usepackage{subcaption}}
\makeatother

\usepackage{bookmark}
\IfFileExists{xurl.sty}{\usepackage{xurl}}{} % add URL line breaks if available
\urlstyle{same}
% fallback for those not using the hyperref driver hyperxmp:
\makeatletter
\@ifundefined{xmpquote}{\newcommand{\xmpquote}[1]{#1}}{}
\makeatother
\hypersetup{
  pdftitle={Africast-Time Series Analysis \& Forecasting Using R},
  hidelinks,
  pdfcreator={LaTeX via pandoc}}


\title{\texorpdfstring{Africast-Time Series Analysis \& Forecasting
Using R}{Africast-Time Series Analysis \& Forecasting Using R}}
\subtitle{\texorpdfstring{4. Transforming / adjusting time
series}{4. Transforming / adjusting time series}}
\author{}
\date{7 October 2026}

\begin{document}
\frame{\titlepage}


\begin{frame}{Outline}
\protect\phantomsection\label{outline}
\vspace*{0.7cm}\tableofcontents
\end{frame}

\section{Transforming time series}\label{transforming-time-series}

\begin{frame}{Variance stabilization}
\protect\phantomsection\label{variance-stabilization}
\fontsize{13}{15}\sf

If the data show different variation at different levels of the series,
then a transformation can be useful. \pause

Denote original observations as \(y_1,\dots,y_n\) and transformed
observations as \(w_1, \dots, w_n\). \pause

\begin{block}{\footnotesize Mathematical transformations for stabilizing
variation}
\begin{tabular}{llc}
Square root & $w_t = \sqrt{y_t}$ & $\downarrow$ \\[0.2cm]
Cube root & $w_t = \sqrt[3]{y_t}$ & Increasing \\[0.2cm]
Logarithm & $w_t = \text{log}(y_t)$  & strength
\end{tabular}
\end{block}
\pause

Logarithms, in particular, are useful because they are more
interpretable: changes in a log value are \textbf{relative (percent)
changes on the original scale}.
\end{frame}

\begin{frame}[fragile]{Variance stabilization}
\protect\phantomsection\label{variance-stabilization-1}
\fontsize{12}{14}\sf

\begin{Shaded}
\begin{Highlighting}[]
\NormalTok{food }\OtherTok{\textless{}{-}}\NormalTok{ aus\_retail }\SpecialCharTok{|\textgreater{}}
  \FunctionTok{filter}\NormalTok{(Industry }\SpecialCharTok{==} \StringTok{"Food retailing"}\NormalTok{) }\SpecialCharTok{|\textgreater{}}
  \FunctionTok{summarise}\NormalTok{(}\AttributeTok{Turnover =} \FunctionTok{sum}\NormalTok{(Turnover))}
\end{Highlighting}
\end{Shaded}

\pandocbounded{\includegraphics[keepaspectratio]{02_transforming_adjusting_files/figure-beamer/food-plot-1.pdf}}
\end{frame}

\begin{frame}[fragile]{Variance stabilization}
\protect\phantomsection\label{variance-stabilization-2}
\fontsize{12}{14}\sf

\begin{Shaded}
\begin{Highlighting}[]
\NormalTok{food }\SpecialCharTok{|\textgreater{}} \FunctionTok{autoplot}\NormalTok{(}\FunctionTok{sqrt}\NormalTok{(Turnover)) }\SpecialCharTok{+}
  \FunctionTok{labs}\NormalTok{(}\AttributeTok{y =} \StringTok{"Square root turnover"}\NormalTok{)}
\end{Highlighting}
\end{Shaded}

\pandocbounded{\includegraphics[keepaspectratio]{02_transforming_adjusting_files/figure-beamer/food-sqrt1-1.pdf}}
\end{frame}

\begin{frame}[fragile]{Variance stabilization}
\protect\phantomsection\label{variance-stabilization-3}
\fontsize{12}{14}\sf

\begin{Shaded}
\begin{Highlighting}[]
\NormalTok{food }\SpecialCharTok{|\textgreater{}} \FunctionTok{autoplot}\NormalTok{(Turnover}\SpecialCharTok{\^{}}\NormalTok{(}\DecValTok{1} \SpecialCharTok{/} \DecValTok{3}\NormalTok{)) }\SpecialCharTok{+}
  \FunctionTok{labs}\NormalTok{(}\AttributeTok{y =} \StringTok{"Cube root turnover"}\NormalTok{)}
\end{Highlighting}
\end{Shaded}

\pandocbounded{\includegraphics[keepaspectratio]{02_transforming_adjusting_files/figure-beamer/food-cbrt-1.pdf}}
\end{frame}

\begin{frame}[fragile]{Variance stabilization}
\protect\phantomsection\label{variance-stabilization-4}
\fontsize{12}{14}\sf

\begin{Shaded}
\begin{Highlighting}[]
\NormalTok{food }\SpecialCharTok{|\textgreater{}} \FunctionTok{autoplot}\NormalTok{(}\FunctionTok{log}\NormalTok{(Turnover)) }\SpecialCharTok{+}
  \FunctionTok{labs}\NormalTok{(}\AttributeTok{y =} \StringTok{"Log turnover"}\NormalTok{)}
\end{Highlighting}
\end{Shaded}

\pandocbounded{\includegraphics[keepaspectratio]{02_transforming_adjusting_files/figure-beamer/food-log-1.pdf}}
\end{frame}

\begin{frame}[fragile]{Variance stabilization}
\protect\phantomsection\label{variance-stabilization-5}
\fontsize{12}{14}\sf

\begin{Shaded}
\begin{Highlighting}[]
\NormalTok{food }\SpecialCharTok{|\textgreater{}} \FunctionTok{autoplot}\NormalTok{(}\SpecialCharTok{{-}}\DecValTok{1} \SpecialCharTok{/}\NormalTok{ Turnover) }\SpecialCharTok{+}
  \FunctionTok{labs}\NormalTok{(}\AttributeTok{y =} \StringTok{"Inverse turnover"}\NormalTok{)}
\end{Highlighting}
\end{Shaded}

\pandocbounded{\includegraphics[keepaspectratio]{02_transforming_adjusting_files/figure-beamer/food-inverse-1.pdf}}
\end{frame}

\begin{frame}{Box-Cox transformations}
\protect\phantomsection\label{box-cox-transformations}
\fontsize{13}{15}\sf

Each of these transformations is close to a member of the family of
\textbf{Box-Cox transformations}: \[w_t = \left\{\begin{array}{ll}
        \text{log}(y_t),      & \quad \lambda = 0; \\
        (\text{sign}(y_t)|y_t|^\lambda-1)/\lambda ,         & \quad \lambda \ne 0.
\end{array}\right.
\]\pause\vspace*{-0.4cm}

\begin{itemize}
\tightlist
\item
  Actually the Bickel-Doksum transformation (allowing for \(y_t<0\))
\item
  \(\lambda=1\): (No substantive transformation)
\item
  \(\lambda=\frac12\): (Square root plus linear transformation)
\item
  \(\lambda=0\): (Natural logarithm)
\item
  \(\lambda=-1\): (Inverse plus 1)
\end{itemize}
\end{frame}

\begin{frame}{Box-Cox transformations}
\protect\phantomsection\label{box-cox-transformations-1}
\animategraphics[,controls,buttonsize=0.3cm,width=11.5cm]{10}{02_transforming_adjusting_files/figure-beamer/food-anim-}{1}{100}
\end{frame}

\begin{frame}[fragile]{Box-Cox transformations}
\protect\phantomsection\label{box-cox-transformations-2}
\fontsize{13}{15}\sf

\begin{Shaded}
\begin{Highlighting}[]
\NormalTok{food }\SpecialCharTok{|\textgreater{}}
  \FunctionTok{features}\NormalTok{(Turnover, }\AttributeTok{features =}\NormalTok{ guerrero)}
\end{Highlighting}
\end{Shaded}

\begin{verbatim}
# A tibble: 1 x 1
  lambda_guerrero
            <dbl>
1          0.0895
\end{verbatim}

\pause

\begin{itemize}
\tightlist
\item
  This attempts to balance the seasonal fluctuations and random
  variation across the series.
\item
  Always check the results.
\item
  A low value of \(\lambda\) can give extremely large prediction
  intervals.
\end{itemize}
\end{frame}

\begin{frame}[fragile]{Box-Cox transformations}
\protect\phantomsection\label{box-cox-transformations-3}
\fontsize{13}{14}\sf

\begin{Shaded}
\begin{Highlighting}[]
\NormalTok{food }\SpecialCharTok{|\textgreater{}} \FunctionTok{autoplot}\NormalTok{(}\FunctionTok{box\_cox}\NormalTok{(Turnover, }\FloatTok{0.0895}\NormalTok{)) }\SpecialCharTok{+}
  \FunctionTok{labs}\NormalTok{(}\AttributeTok{y =} \StringTok{"Box{-}Cox transformed turnover"}\NormalTok{)}
\end{Highlighting}
\end{Shaded}

\pandocbounded{\includegraphics[keepaspectratio]{02_transforming_adjusting_files/figure-beamer/food-bc-1.pdf}}
\end{frame}

\begin{frame}[fragile]{Transformations}
\protect\phantomsection\label{transformations}
\fontsize{13}{15}\sf

\begin{itemize}
\tightlist
\item
  Often no transformation needed.
\item
  Simple transformations are easier to explain and work well enough.
\item
  Transformations can have very large effect on PI.
\item
  If some data are zero or negative, then use \(\lambda>0\).
\item
  \texttt{log1p()} can also be useful for data with zeros.
\item
  Choosing logs is a simple way to force forecasts to be positive
\item
  Transformations must be reversed to obtain forecasts on the original
  scale. (Handled automatically by \texttt{fable}.)
\end{itemize}
\end{frame}

\section{Adjusting time series}\label{adjusting-time-series}

\begin{frame}[fragile]{Per capita adjustments}
\protect\phantomsection\label{per-capita-adjustments}
\begin{Shaded}
\begin{Highlighting}[]
\NormalTok{global\_economy }\SpecialCharTok{|\textgreater{}}
  \FunctionTok{filter}\NormalTok{(Country }\SpecialCharTok{==} \StringTok{"Australia"}\NormalTok{) }\SpecialCharTok{|\textgreater{}}
  \FunctionTok{autoplot}\NormalTok{(GDP)}
\end{Highlighting}
\end{Shaded}

\pandocbounded{\includegraphics[keepaspectratio]{02_transforming_adjusting_files/figure-beamer/gdp-per-capita-1.pdf}}
\end{frame}

\begin{frame}[fragile]{Per capita adjustments}
\protect\phantomsection\label{per-capita-adjustments-1}
\begin{Shaded}
\begin{Highlighting}[]
\NormalTok{global\_economy }\SpecialCharTok{|\textgreater{}}
  \FunctionTok{filter}\NormalTok{(Country }\SpecialCharTok{==} \StringTok{"Australia"}\NormalTok{) }\SpecialCharTok{|\textgreater{}}
  \FunctionTok{autoplot}\NormalTok{(GDP }\SpecialCharTok{/}\NormalTok{ Population)}
\end{Highlighting}
\end{Shaded}

\pandocbounded{\includegraphics[keepaspectratio]{02_transforming_adjusting_files/figure-beamer/gdp-per-capita2-1.pdf}}
\end{frame}

\begin{frame}[fragile]{Inflation adjustments}
\protect\phantomsection\label{inflation-adjustments}
\fontsize{10}{10}\sf

\begin{Shaded}
\begin{Highlighting}[]
\NormalTok{print\_retail }\OtherTok{\textless{}{-}}\NormalTok{ aus\_retail }\SpecialCharTok{|\textgreater{}}
  \FunctionTok{filter}\NormalTok{(Industry }\SpecialCharTok{==} \StringTok{"Newspaper and book retailing"}\NormalTok{) }\SpecialCharTok{|\textgreater{}}
  \FunctionTok{group\_by}\NormalTok{(Industry) }\SpecialCharTok{|\textgreater{}}
  \FunctionTok{index\_by}\NormalTok{(}\AttributeTok{Year =} \FunctionTok{year}\NormalTok{(Month)) }\SpecialCharTok{|\textgreater{}}
  \FunctionTok{summarise}\NormalTok{(}\AttributeTok{Turnover =} \FunctionTok{sum}\NormalTok{(Turnover))}
\NormalTok{aus\_economy }\OtherTok{\textless{}{-}} \FunctionTok{filter}\NormalTok{(global\_economy, Country }\SpecialCharTok{==} \StringTok{"Australia"}\NormalTok{)}
\NormalTok{print\_retail }\SpecialCharTok{|\textgreater{}}
  \FunctionTok{left\_join}\NormalTok{(aus\_economy, }\AttributeTok{by =} \StringTok{"Year"}\NormalTok{) }\SpecialCharTok{|\textgreater{}}
  \FunctionTok{mutate}\NormalTok{(}\AttributeTok{Adj\_turnover =}\NormalTok{ Turnover }\SpecialCharTok{/}\NormalTok{ CPI) }\SpecialCharTok{|\textgreater{}}
  \FunctionTok{pivot\_longer}\NormalTok{(}\FunctionTok{c}\NormalTok{(Turnover, Adj\_turnover),}
    \AttributeTok{names\_to =} \StringTok{"Type"}\NormalTok{, }\AttributeTok{values\_to =} \StringTok{"Turnover"}
\NormalTok{  ) }\SpecialCharTok{|\textgreater{}}
  \FunctionTok{ggplot}\NormalTok{(}\FunctionTok{aes}\NormalTok{(}\AttributeTok{x =}\NormalTok{ Year, }\AttributeTok{y =}\NormalTok{ Turnover)) }\SpecialCharTok{+}
  \FunctionTok{geom\_line}\NormalTok{() }\SpecialCharTok{+}
  \FunctionTok{facet\_grid}\NormalTok{(}\FunctionTok{vars}\NormalTok{(Type), }\AttributeTok{scales =} \StringTok{"free\_y"}\NormalTok{) }\SpecialCharTok{+}
  \FunctionTok{labs}\NormalTok{(}\AttributeTok{x =} \StringTok{"Years"}\NormalTok{, }\AttributeTok{y =} \ConstantTok{NULL}\NormalTok{,}
       \AttributeTok{title =} \StringTok{"Turnover: Australian print media industry"}\NormalTok{)}
\end{Highlighting}
\end{Shaded}
\end{frame}

\begin{frame}{Inflation adjustments}
\protect\phantomsection\label{inflation-adjustments-1}
\fontsize{10}{10}\sf

\includegraphics[width=\linewidth,height=0.9\textheight,keepaspectratio]{02_transforming_adjusting_files/figure-beamer/unnamed-chunk-2-1.pdf}
\end{frame}

\section{Time series decompositions}\label{time-series-decompositions}

\begin{frame}{Time series decomposition}
\protect\phantomsection\label{time-series-decomposition}
\begin{description}
\tightlist
\item[Trend-Cycle]
aperiodic changes in level over time.
\item[Seasonal]
(almost) periodic changes in level due to seasonal factors (e.g., the
quarter of the year, the month, or day of the week).
\end{description}

\begin{block}{Additive decomposition}\vspace*{-0.3cm}
\[ y_t = S_t + T_t + R_t \]
\end{block}
\begin{tabular}{@{}llp{8cm}@{}}
where & $y_t=$ & data at period $t$ \\
      & $T_t=$ & trend-cycle component at period $t$\\
      & $S_t=$ & seasonal component at period $t$ \\
      & $R_t=$ & remainder component at period $t$
\end{tabular}
\end{frame}

\begin{frame}{STL decomposition}
\protect\phantomsection\label{stl-decomposition}
\fontsize{13}{14}\sf

\begin{itemize}
\tightlist
\item
  STL: ``Seasonal and Trend decomposition using Loess''
\item
  Very versatile and robust.
\item
  Seasonal component allowed to change over time, and rate of change
  controlled by user.
\item
  Smoothness of trend-cycle also controlled by user.
\item
  Optionally robust to outliers
\item
  No trading day or calendar adjustments.
\item
  Only additive.
\item
  Take logs to get multiplicative decomposition.
\item
  Use Box-Cox transformations to get other decompositions.
\end{itemize}
\end{frame}

\begin{frame}[fragile]{US Retail Employment}
\protect\phantomsection\label{us-retail-employment}
\fontsize{11}{11}\sf

\begin{Shaded}
\begin{Highlighting}[]
\NormalTok{us\_retail\_employment }\OtherTok{\textless{}{-}}\NormalTok{ us\_employment }\SpecialCharTok{|\textgreater{}}
  \FunctionTok{filter}\NormalTok{(}\FunctionTok{year}\NormalTok{(Month) }\SpecialCharTok{\textgreater{}=} \DecValTok{1990}\NormalTok{, Title }\SpecialCharTok{==} \StringTok{"Retail Trade"}\NormalTok{) }\SpecialCharTok{|\textgreater{}}
  \FunctionTok{select}\NormalTok{(}\SpecialCharTok{{-}}\NormalTok{Series\_ID)}
\NormalTok{us\_retail\_employment}
\end{Highlighting}
\end{Shaded}

\begin{verbatim}
# A tsibble: 357 x 3 [1M]
      Month Title        Employed
      <mth> <chr>           <dbl>
 1 1990 Jan Retail Trade   13256.
 2 1990 Feb Retail Trade   12966.
 3 1990 Mar Retail Trade   12938.
 4 1990 Apr Retail Trade   13012.
 5 1990 May Retail Trade   13108.
 6 1990 Jun Retail Trade   13183.
 7 1990 Jul Retail Trade   13170.
 8 1990 Aug Retail Trade   13160.
 9 1990 Sep Retail Trade   13113.
10 1990 Oct Retail Trade   13185.
# i 347 more rows
\end{verbatim}

\vspace*{10cm}
\end{frame}

\begin{frame}[fragile]{US Retail Employment}
\protect\phantomsection\label{us-retail-employment-1}
\begin{Shaded}
\begin{Highlighting}[]
\NormalTok{us\_retail\_employment }\SpecialCharTok{|\textgreater{}}
  \FunctionTok{autoplot}\NormalTok{(Employed) }\SpecialCharTok{+}
  \FunctionTok{labs}\NormalTok{(}\AttributeTok{y =} \StringTok{"Persons (thousands)"}\NormalTok{, }\AttributeTok{title =} \StringTok{"Total employment in US retail"}\NormalTok{)}
\end{Highlighting}
\end{Shaded}

\pandocbounded{\includegraphics[keepaspectratio]{02_transforming_adjusting_files/figure-beamer/dable1-1.pdf}}

\vspace*{10cm}
\end{frame}

\begin{frame}[fragile]{US Retail Employment}
\protect\phantomsection\label{us-retail-employment-2}
\fontsize{11}{11}\sf

\begin{Shaded}
\begin{Highlighting}[]
\NormalTok{dcmp }\OtherTok{\textless{}{-}}\NormalTok{ us\_retail\_employment }\SpecialCharTok{|\textgreater{}}
  \FunctionTok{model}\NormalTok{(}\AttributeTok{stl =} \FunctionTok{STL}\NormalTok{(Employed))}
\NormalTok{dcmp}
\end{Highlighting}
\end{Shaded}

\begin{verbatim}
# A mable: 1 x 1
      stl
  <model>
1   <STL>
\end{verbatim}

\vspace*{10cm}
\end{frame}

\begin{frame}[fragile]{US Retail Employment}
\protect\phantomsection\label{us-retail-employment-3}
\fontsize{11}{11}\sf

\begin{Shaded}
\begin{Highlighting}[]
\FunctionTok{components}\NormalTok{(dcmp)}
\end{Highlighting}
\end{Shaded}

\begin{verbatim}
# A dable: 357 x 7 [1M]
# Key:     .model [1]
# :        Employed = trend + season_year + remainder
   .model    Month Employed  trend season_year remainder season_adjust
   <chr>     <mth>    <dbl>  <dbl>       <dbl>     <dbl>         <dbl>
 1 stl    1990 Jan   13256. 13288.      -33.0      0.836        13289.
 2 stl    1990 Feb   12966. 13269.     -258.     -44.6          13224.
 3 stl    1990 Mar   12938. 13250.     -290.     -22.1          13228.
 4 stl    1990 Apr   13012. 13231.     -220.       1.05         13232.
 5 stl    1990 May   13108. 13211.     -114.      11.3          13223.
 6 stl    1990 Jun   13183. 13192.      -24.3     15.5          13207.
 7 stl    1990 Jul   13170. 13172.      -23.2     21.6          13193.
 8 stl    1990 Aug   13160. 13151.       -9.52    17.8          13169.
 9 stl    1990 Sep   13113. 13131.      -39.5     22.0          13153.
10 stl    1990 Oct   13185. 13110.       61.6     13.2          13124.
# i 347 more rows
\end{verbatim}

\vspace*{10cm}
\end{frame}

\begin{frame}[fragile]{US Retail Employment}
\protect\phantomsection\label{us-retail-employment-4}
\begin{Shaded}
\begin{Highlighting}[]
\FunctionTok{components}\NormalTok{(dcmp) }\SpecialCharTok{|\textgreater{}} \FunctionTok{autoplot}\NormalTok{()}
\end{Highlighting}
\end{Shaded}

\pandocbounded{\includegraphics[keepaspectratio]{02_transforming_adjusting_files/figure-beamer/usretail-stl-1.pdf}}
\end{frame}

\begin{frame}[fragile]{US Retail Employment}
\protect\phantomsection\label{us-retail-employment-5}
\begin{Shaded}
\begin{Highlighting}[]
\NormalTok{us\_retail\_employment }\SpecialCharTok{|\textgreater{}}
  \FunctionTok{autoplot}\NormalTok{(Employed, }\AttributeTok{color =} \StringTok{"gray"}\NormalTok{) }\SpecialCharTok{+}
  \FunctionTok{autolayer}\NormalTok{(}\FunctionTok{components}\NormalTok{(dcmp), trend, }\AttributeTok{color =} \StringTok{"\#D55E00"}\NormalTok{) }\SpecialCharTok{+}
  \FunctionTok{labs}\NormalTok{(}\AttributeTok{y =} \StringTok{"Persons (thousands)"}\NormalTok{, }\AttributeTok{title =} \StringTok{"Total employment in US retail"}\NormalTok{)}
\end{Highlighting}
\end{Shaded}

\pandocbounded{\includegraphics[keepaspectratio]{02_transforming_adjusting_files/figure-beamer/dable4-1.pdf}}

\vspace*{10cm}
\end{frame}

\begin{frame}[fragile]{US Retail Employment}
\protect\phantomsection\label{us-retail-employment-6}
\begin{Shaded}
\begin{Highlighting}[]
\FunctionTok{components}\NormalTok{(dcmp) }\SpecialCharTok{|\textgreater{}} \FunctionTok{gg\_subseries}\NormalTok{(season\_year)}
\end{Highlighting}
\end{Shaded}

\pandocbounded{\includegraphics[keepaspectratio]{02_transforming_adjusting_files/figure-beamer/usretail3-1.pdf}}
\end{frame}

\begin{frame}{Seasonal adjustment}
\protect\phantomsection\label{seasonal-adjustment}
\begin{itemize}
\tightlist
\item
  Useful by-product of decomposition: an easy way to calculate
  seasonally adjusted data.
\item
  Additive decomposition: seasonally adjusted data given by
  \[y_t - S_t = T_t + R_t\]
\item
  Multiplicative decomposition: seasonally adjusted data given by
  \[y_t / S_t = T_t \times R_t\]
\end{itemize}
\end{frame}

\begin{frame}[fragile]{US Retail Employment}
\protect\phantomsection\label{us-retail-employment-7}
\begin{Shaded}
\begin{Highlighting}[]
\NormalTok{us\_retail\_employment }\SpecialCharTok{|\textgreater{}}
  \FunctionTok{autoplot}\NormalTok{(Employed, }\AttributeTok{color =} \StringTok{"gray"}\NormalTok{) }\SpecialCharTok{+}
  \FunctionTok{autolayer}\NormalTok{(}\FunctionTok{components}\NormalTok{(dcmp), season\_adjust, }\AttributeTok{color =} \StringTok{"\#0072B2"}\NormalTok{) }\SpecialCharTok{+}
  \FunctionTok{labs}\NormalTok{(}\AttributeTok{y =} \StringTok{"Persons (thousands)"}\NormalTok{, }\AttributeTok{title =} \StringTok{"Total employment in US retail"}\NormalTok{)}
\end{Highlighting}
\end{Shaded}

\pandocbounded{\includegraphics[keepaspectratio]{02_transforming_adjusting_files/figure-beamer/usretail-sa-1.pdf}}
\end{frame}

\begin{frame}{Seasonal adjustment}
\protect\phantomsection\label{seasonal-adjustment-1}
\begin{itemize}
\tightlist
\item
  We use estimates of \(S\) based on past values to seasonally adjust a
  current value.
\item
  Seasonally adjusted series reflect \textbf{remainders} as well as
  \textbf{trend}. Therefore they are not ``smooth'' and ``downturns'' or
  ``upturns'' can be misleading.
\item
  It is better to use the trend-cycle component to look for turning
  points.
\end{itemize}
\end{frame}

\begin{frame}{STL decomposition}
\protect\phantomsection\label{stl-decomposition-1}
\animategraphics[,controls,buttonsize=0.3cm,width=11.5cm]{10}{02_transforming_adjusting_files/figure-beamer/stlwindowanim-}{1}{100}

\vspace*{10cm}
\end{frame}

\begin{frame}[fragile]{STL decomposition}
\protect\phantomsection\label{stl-decomposition-2}
\begin{Shaded}
\begin{Highlighting}[]
\NormalTok{us\_retail\_employment }\SpecialCharTok{|\textgreater{}}
  \FunctionTok{model}\NormalTok{(}\FunctionTok{STL}\NormalTok{(Employed }\SpecialCharTok{\textasciitilde{}} \FunctionTok{trend}\NormalTok{(}\AttributeTok{window =} \DecValTok{15}\NormalTok{) }\SpecialCharTok{+} \FunctionTok{season}\NormalTok{(}\AttributeTok{window =} \StringTok{"periodic"}\NormalTok{),}
    \AttributeTok{robust =} \ConstantTok{TRUE}
\NormalTok{  )) }\SpecialCharTok{|\textgreater{}}
  \FunctionTok{components}\NormalTok{()}
\end{Highlighting}
\end{Shaded}

\fontsize{12}{13}\sf

\begin{itemize}
\tightlist
\item
  \texttt{trend(window\ =\ ?)} controls wiggliness of trend component.
\item
  \texttt{season(window\ =\ ?)} controls variation on seasonal
  component.
\item
  \texttt{season(window\ =\ \textquotesingle{}periodic\textquotesingle{})}
  is equivalent to an infinite window.
\end{itemize}
\end{frame}

\begin{frame}[fragile]{STL decomposition}
\protect\phantomsection\label{stl-decomposition-3}
\begin{Shaded}
\begin{Highlighting}[]
\NormalTok{us\_retail\_employment }\SpecialCharTok{|\textgreater{}}
  \FunctionTok{model}\NormalTok{(}\FunctionTok{STL}\NormalTok{(Employed)) }\SpecialCharTok{|\textgreater{}}
  \FunctionTok{components}\NormalTok{() }\SpecialCharTok{|\textgreater{}}
  \FunctionTok{autoplot}\NormalTok{()}
\end{Highlighting}
\end{Shaded}

\pandocbounded{\includegraphics[keepaspectratio]{02_transforming_adjusting_files/figure-beamer/mstl-1.pdf}}

\only<2>{\begin{textblock}{7}(8,0.2)\fontsize{11}{11}\sf
\begin{alertblock}{}
\begin{itemize}\tightlist
\item \texttt{STL()} chooses \texttt{season(window=13)} by default
\item Can include transformations.
\end{itemize}
\end{alertblock}
\end{textblock}}
\end{frame}

\begin{frame}[fragile]{STL decomposition}
\protect\phantomsection\label{stl-decomposition-4}
\fontsize{13}{15}\sf

\begin{itemize}
\tightlist
\item
  Algorithm that updates trend and seasonal components iteratively.
\item
  Starts with \(\hat{T}_t=0\)
\item
  Uses a mixture of loess and moving averages to successively refine the
  trend and seasonal estimates.
\item
  trend window controls loess bandwidth on deasonalised values.
\item
  season window controls loess bandwidth on detrended subseries.
\item
  Robustness weights based on remainder.
\item
  Default season: \texttt{window\ =\ 13}
\item
  Default
  trend:\newline\mbox{}\hfill \texttt{window\ =\ nextodd(ceiling((1.5*period)/(1-(1.5/s.window)))}
\item
  \texttt{window} values should be odd numbers for symmetry.
\end{itemize}
\end{frame}

\begin{frame}[fragile]{Australian holidays}
\protect\phantomsection\label{australian-holidays}
\begin{Shaded}
\begin{Highlighting}[]
\NormalTok{holidays }\SpecialCharTok{|\textgreater{}} \FunctionTok{autoplot}\NormalTok{(Trips) }\SpecialCharTok{+}
  \FunctionTok{labs}\NormalTok{(}\AttributeTok{y =} \StringTok{"thousands of trips"}\NormalTok{, }\AttributeTok{x =} \StringTok{"Year"}\NormalTok{,}
       \AttributeTok{title =} \StringTok{"Australian domestic holiday nights"}\NormalTok{)}
\end{Highlighting}
\end{Shaded}

\pandocbounded{\includegraphics[keepaspectratio]{02_transforming_adjusting_files/figure-beamer/holidays-plot2-1.pdf}}
\end{frame}

\begin{frame}[fragile]{Australian holidays}
\protect\phantomsection\label{australian-holidays-1}
\begin{Shaded}
\begin{Highlighting}[]
\NormalTok{holidays }\SpecialCharTok{|\textgreater{}}
  \FunctionTok{model}\NormalTok{(}\AttributeTok{stl =} \FunctionTok{STL}\NormalTok{(Trips)) }\SpecialCharTok{|\textgreater{}}
  \FunctionTok{components}\NormalTok{() }\SpecialCharTok{|\textgreater{}}
  \FunctionTok{autoplot}\NormalTok{()}
\end{Highlighting}
\end{Shaded}

\pandocbounded{\includegraphics[keepaspectratio]{02_transforming_adjusting_files/figure-beamer/stlagain2-1.pdf}}
\end{frame}

\begin{frame}[fragile]{Holidays decomposition}
\protect\phantomsection\label{holidays-decomposition}
\fontsize{9}{10}\sf

\begin{Shaded}
\begin{Highlighting}[]
\NormalTok{dcmp }\OtherTok{\textless{}{-}}\NormalTok{ holidays }\SpecialCharTok{|\textgreater{}}
  \FunctionTok{model}\NormalTok{(}\AttributeTok{stl =} \FunctionTok{STL}\NormalTok{(Trips)) }\SpecialCharTok{|\textgreater{}}
  \FunctionTok{components}\NormalTok{()}
\NormalTok{dcmp}
\end{Highlighting}
\end{Shaded}

\begin{verbatim}
# A dable: 640 x 8 [1Q]
# Key:     State, .model [8]
# :        Trips = trend + season_year + remainder
   State .model Quarter Trips trend season_year remainder season_adjust
   <chr> <chr>    <qtr> <dbl> <dbl>       <dbl>     <dbl>         <dbl>
 1 ACT   stl    1998 Q1  196.  172.       -8.48     32.6           205.
 2 ACT   stl    1998 Q2  127.  157.       10.3     -40.6           116.
 3 ACT   stl    1998 Q3  111.  142.      -16.8     -14.5           128.
 4 ACT   stl    1998 Q4  170.  130.       14.6      25.6           156.
 5 ACT   stl    1999 Q1  108.  135.       -8.63    -18.3           116.
 6 ACT   stl    1999 Q2  125.  148.       11.0     -34.6           114.
 7 ACT   stl    1999 Q3  178.  166.      -16.0      28.3           194.
 8 ACT   stl    1999 Q4  218.  177.       13.2      27.5           204.
 9 ACT   stl    2000 Q1  158.  169.       -8.75     -1.96          167.
10 ACT   stl    2000 Q2  155.  151.       11.7      -8.20          143.
# i 630 more rows
\end{verbatim}
\end{frame}

\begin{frame}[fragile]{Holidays decomposition}
\protect\phantomsection\label{holidays-decomposition-1}
\fontsize{9}{10}\sf

\begin{Shaded}
\begin{Highlighting}[]
\NormalTok{dcmp }\SpecialCharTok{|\textgreater{}} \FunctionTok{gg\_subseries}\NormalTok{(season\_year)}
\end{Highlighting}
\end{Shaded}

\pandocbounded{\includegraphics[keepaspectratio]{02_transforming_adjusting_files/figure-beamer/holidays3-1.pdf}}
\end{frame}

\begin{frame}[fragile]{Holidays decomposition}
\protect\phantomsection\label{holidays-decomposition-2}
\fontsize{9}{10}\sf

\begin{Shaded}
\begin{Highlighting}[]
\FunctionTok{autoplot}\NormalTok{(dcmp, trend, }\AttributeTok{scale\_bars =} \ConstantTok{FALSE}\NormalTok{) }\SpecialCharTok{+}
  \FunctionTok{autolayer}\NormalTok{(holidays, }\AttributeTok{alpha =} \FloatTok{0.4}\NormalTok{)}
\end{Highlighting}
\end{Shaded}

\includegraphics[width=\linewidth,height=0.7\textheight,keepaspectratio]{02_transforming_adjusting_files/figure-beamer/holidays-trend-1.pdf}
\end{frame}

\section{Multiple seasonality}\label{multiple-seasonality}

\begin{frame}[fragile]{Multiple seasonality}
\protect\phantomsection\label{multiple-seasonality-1}
\fontsize{9}{9}\sf

\begin{Shaded}
\begin{Highlighting}[]
\NormalTok{vic\_elec }\SpecialCharTok{|\textgreater{}}
  \FunctionTok{model}\NormalTok{(}\FunctionTok{STL}\NormalTok{(Demand)) }\SpecialCharTok{|\textgreater{}}
  \FunctionTok{components}\NormalTok{() }\SpecialCharTok{|\textgreater{}}
  \FunctionTok{autoplot}\NormalTok{()}
\end{Highlighting}
\end{Shaded}

\placefig{7}{0.2}{width=8.7cm, height=20cm}{vic_elec_stl}

\vspace*{10cm}
\end{frame}

\section{The ABS stuff-up}\label{the-abs-stuff-up}

\begin{frame}{The ABS stuff-up}
\protect\phantomsection\label{the-abs-stuff-up-1}
\full{abs1}
\end{frame}

\begin{frame}{The ABS stuff-up}
\protect\phantomsection\label{the-abs-stuff-up-2}
\full{abs2}
\end{frame}

\begin{frame}{The ABS stuff-up}
\protect\phantomsection\label{the-abs-stuff-up-3}
\full{abs3}
\end{frame}

\begin{frame}[fragile]{The ABS stuff-up}
\protect\phantomsection\label{the-abs-stuff-up-4}
\fontsize{11}{11}\sf

\begin{Shaded}
\begin{Highlighting}[]
\NormalTok{employed}
\end{Highlighting}
\end{Shaded}

\begin{verbatim}
# A tsibble: 440 x 4 [1M]
       Time Month  Year Employed
      <mth> <ord> <dbl>    <dbl>
 1 1978 Feb Feb    1978    5986.
 2 1978 Mar Mar    1978    6041.
 3 1978 Apr Apr    1978    6054.
 4 1978 May May    1978    6038.
 5 1978 Jun Jun    1978    6031.
 6 1978 Jul Jul    1978    6036.
 7 1978 Aug Aug    1978    6005.
 8 1978 Sep Sep    1978    6024.
 9 1978 Oct Oct    1978    6046.
10 1978 Nov Nov    1978    6034.
# i 430 more rows
\end{verbatim}
\end{frame}

\begin{frame}[fragile]{The ABS stuff-up}
\protect\phantomsection\label{the-abs-stuff-up-5}
\begin{Shaded}
\begin{Highlighting}[]
\NormalTok{employed }\SpecialCharTok{|\textgreater{}}
  \FunctionTok{autoplot}\NormalTok{(Employed) }\SpecialCharTok{+}
  \FunctionTok{labs}\NormalTok{(}\AttributeTok{title =} \StringTok{"Total employed"}\NormalTok{, }\AttributeTok{y =} \StringTok{"Thousands"}\NormalTok{)}
\end{Highlighting}
\end{Shaded}

\pandocbounded{\includegraphics[keepaspectratio]{02_transforming_adjusting_files/figure-beamer/abs3-1.pdf}}
\end{frame}

\begin{frame}[fragile]{The ABS stuff-up}
\protect\phantomsection\label{the-abs-stuff-up-6}
\begin{Shaded}
\begin{Highlighting}[]
\NormalTok{employed }\SpecialCharTok{|\textgreater{}}
  \FunctionTok{filter}\NormalTok{(Year }\SpecialCharTok{\textgreater{}=} \DecValTok{2005}\NormalTok{) }\SpecialCharTok{|\textgreater{}}
  \FunctionTok{autoplot}\NormalTok{(Employed) }\SpecialCharTok{+}
  \FunctionTok{labs}\NormalTok{(}\AttributeTok{title =} \StringTok{"Total employed"}\NormalTok{, }\AttributeTok{y =} \StringTok{"Thousands"}\NormalTok{)}
\end{Highlighting}
\end{Shaded}

\pandocbounded{\includegraphics[keepaspectratio]{02_transforming_adjusting_files/figure-beamer/abs4-1.pdf}}
\end{frame}

\begin{frame}[fragile]{The ABS stuff-up}
\protect\phantomsection\label{the-abs-stuff-up-7}
\begin{Shaded}
\begin{Highlighting}[]
\NormalTok{employed }\SpecialCharTok{|\textgreater{}}
  \FunctionTok{filter}\NormalTok{(Year }\SpecialCharTok{\textgreater{}=} \DecValTok{2005}\NormalTok{) }\SpecialCharTok{|\textgreater{}}
  \FunctionTok{gg\_season}\NormalTok{(Employed, }\AttributeTok{labels =} \StringTok{"right"}\NormalTok{) }\SpecialCharTok{+}
  \FunctionTok{labs}\NormalTok{(}\AttributeTok{title =} \StringTok{"Total employed"}\NormalTok{, }\AttributeTok{y =} \StringTok{"Thousands"}\NormalTok{)}
\end{Highlighting}
\end{Shaded}

\pandocbounded{\includegraphics[keepaspectratio]{02_transforming_adjusting_files/figure-beamer/abs5-1.pdf}}
\end{frame}

\begin{frame}[fragile]{The ABS stuff-up}
\protect\phantomsection\label{the-abs-stuff-up-8}
\begin{Shaded}
\begin{Highlighting}[]
\NormalTok{employed }\SpecialCharTok{|\textgreater{}}
  \FunctionTok{mutate}\NormalTok{(}\AttributeTok{diff =} \FunctionTok{difference}\NormalTok{(Employed)) }\SpecialCharTok{|\textgreater{}}
  \FunctionTok{filter}\NormalTok{(Month }\SpecialCharTok{==} \StringTok{"Sep"}\NormalTok{) }\SpecialCharTok{|\textgreater{}}
  \FunctionTok{ggplot}\NormalTok{(}\FunctionTok{aes}\NormalTok{(}\AttributeTok{y =}\NormalTok{ diff, }\AttributeTok{x =} \DecValTok{1}\NormalTok{)) }\SpecialCharTok{+}
  \FunctionTok{geom\_boxplot}\NormalTok{() }\SpecialCharTok{+}
  \FunctionTok{coord\_flip}\NormalTok{() }\SpecialCharTok{+}
  \FunctionTok{labs}\NormalTok{(}\AttributeTok{title =} \StringTok{"Sep {-} Aug: total employed"}\NormalTok{, }\AttributeTok{y =} \StringTok{"Thousands"}\NormalTok{) }\SpecialCharTok{+}
  \FunctionTok{scale\_x\_continuous}\NormalTok{(}\AttributeTok{breaks =} \ConstantTok{NULL}\NormalTok{, }\AttributeTok{labels =} \ConstantTok{NULL}\NormalTok{)}
\end{Highlighting}
\end{Shaded}

\pandocbounded{\includegraphics[keepaspectratio]{02_transforming_adjusting_files/figure-beamer/abs6-1.pdf}}
\end{frame}

\begin{frame}[fragile]{The ABS stuff-up}
\protect\phantomsection\label{the-abs-stuff-up-9}
\begin{Shaded}
\begin{Highlighting}[]
\NormalTok{dcmp }\OtherTok{\textless{}{-}}\NormalTok{ employed }\SpecialCharTok{|\textgreater{}}
  \FunctionTok{filter}\NormalTok{(Year }\SpecialCharTok{\textgreater{}=} \DecValTok{2005}\NormalTok{) }\SpecialCharTok{|\textgreater{}}
  \FunctionTok{model}\NormalTok{(}\AttributeTok{stl =} \FunctionTok{STL}\NormalTok{(Employed }\SpecialCharTok{\textasciitilde{}} \FunctionTok{season}\NormalTok{(}\AttributeTok{window =} \DecValTok{11}\NormalTok{), }\AttributeTok{robust =} \ConstantTok{TRUE}\NormalTok{))}
\FunctionTok{components}\NormalTok{(dcmp) }\SpecialCharTok{|\textgreater{}} \FunctionTok{autoplot}\NormalTok{()}
\end{Highlighting}
\end{Shaded}

\pandocbounded{\includegraphics[keepaspectratio]{02_transforming_adjusting_files/figure-beamer/abs7-1.pdf}}
\end{frame}

\begin{frame}[fragile]{The ABS stuff-up}
\protect\phantomsection\label{the-abs-stuff-up-10}
\begin{Shaded}
\begin{Highlighting}[]
\FunctionTok{components}\NormalTok{(dcmp) }\SpecialCharTok{|\textgreater{}}
  \FunctionTok{filter}\NormalTok{(}\FunctionTok{year}\NormalTok{(Time) }\SpecialCharTok{==} \DecValTok{2013}\NormalTok{) }\SpecialCharTok{|\textgreater{}}
  \FunctionTok{gg\_season}\NormalTok{(season\_year) }\SpecialCharTok{+}
  \FunctionTok{labs}\NormalTok{(}\AttributeTok{title =} \StringTok{"Seasonal component"}\NormalTok{) }\SpecialCharTok{+} \FunctionTok{guides}\NormalTok{(}\AttributeTok{colour =} \StringTok{"none"}\NormalTok{)}
\end{Highlighting}
\end{Shaded}

\pandocbounded{\includegraphics[keepaspectratio]{02_transforming_adjusting_files/figure-beamer/abs8-1.pdf}}
\end{frame}

\begin{frame}[fragile]{The ABS stuff-up}
\protect\phantomsection\label{the-abs-stuff-up-11}
\begin{Shaded}
\begin{Highlighting}[]
\FunctionTok{components}\NormalTok{(dcmp) }\SpecialCharTok{|\textgreater{}}
  \FunctionTok{as\_tsibble}\NormalTok{() }\SpecialCharTok{|\textgreater{}}
  \FunctionTok{autoplot}\NormalTok{(season\_adjust)}
\end{Highlighting}
\end{Shaded}

\pandocbounded{\includegraphics[keepaspectratio]{02_transforming_adjusting_files/figure-beamer/abs9-1.pdf}}
\end{frame}

\begin{frame}{The ABS stuff-up}
\protect\phantomsection\label{the-abs-stuff-up-12}
\fontsize{13}{15}\sf

\begin{itemize}
\tightlist
\item
  August 2014 employment numbers higher than expected.
\item
  Supplementary survey usually conducted in August for employed people.
\item
  Most likely, some employed people were claiming to be unemployed in
  August to avoid supplementary questions.
\item
  Supplementary survey not run in 2014, so no motivation to lie about
  employment.
\item
  In previous years, seasonal adjustment fixed the problem.
\item
  The ABS has now adopted a new method to avoid the bias.
\end{itemize}
\end{frame}




\end{document}
